% Seccion: Caso C: Fourier 2D y ruido periódico
% Responsable: R1 (DFT 2D, fase) y R2 (ruido periódico, detección). Extension: 1,5 pagina(s).
% Solo R1 (DFT 2D, fase) y R2 (ruido periódico, detección) edita este archivo.
% R1 escribe la subseccion "DFT 2D y fase"; R2 escribe "Ruido periódico y detección de picos".

\section{Caso C: Fourier 2D y ruido periódico}
\label{sec:fourier2d}

\subsection{DFT 2D y fase}
\label{sec:dft2}

\textit{(Por escribir, R1.)}

\subsection{Ruido periódico y detección de picos}
\label{sec:ruido}

Modelamos la interferencia periódica como un coseno 2D sumado a la imagen,
\begin{equation}
  r(x,y) = A\cos\!\left(2\pi\left(\frac{k_u\,x}{M} + \frac{k_v\,y}{N}\right)\right),
  \label{eq:coseno2d}
\end{equation}
con $x$ la columna, $y$ la fila e imagen de $N\times M$ píxeles. Elegimos $k_u$ y $k_v$ enteros,
es decir, un número exacto de ciclos en toda la imagen. Así la frecuencia cae justo en un bin de
la DFT 2D y aparecen exactamente dos picos, en $(k_v, k_u)$ y en su conjugado
$(N-k_v,\,M-k_u)$, cada uno de altura $A\,NM/2$. Si los ciclos no fueran enteros habría fuga, como en la
Sección~\ref{sec:dft}. Una prueba unitaria verifica que el coseno 2D da exactamente esos dos
picos.

Para encontrar los picos de forma automática (\texttt{detectar\_picos}), tomamos $|E|$, ponemos en
cero un disco de radio $r_{\text{DC}}$ alrededor del DC, nos quedamos con los bins que superan una
fracción \texttt{umbral} del máximo que queda, y agregamos siempre el conjugado de cada pico.
Esto último es necesario: si el notch borrara solo un pico del par, la imagen filtrada dejaría
de ser real.

El disco alrededor del DC es lo que hace funcionar esta versión simple. Una imagen natural tiene
casi toda su energía en frecuencias bajas, y sin el disco esos bins le ganan a cualquier
interferencia débil. Para elegir $r_{\text{DC}}$ probamos con la imagen \texttt{camera} y dos
interferencias de amplitud decreciente (Tabla~\ref{tab:sensibilidad}). Con $r_{\text{DC}}=10$
el detector encuentra los 4 picos sin ningún falso hasta una amplitud de 0{,}05; con
$r_{\text{DC}}=4$ ya aparecen falsos con 0{,}1. Bajo 0{,}02 ninguno de los dos funciona bien, y
ese es el límite del método. El costo es no detectar patrones de menos de 10 ciclos en la imagen,
que de todas formas se confunden con el contenido. Usamos \texttt{umbral} $=0{,}3$.

\begin{table}[H]
  \centering
  \small
  \begin{tabular}{ccccc}
    \toprule
    Amplitud & \multicolumn{2}{c}{$r_{\text{DC}}=4$} & \multicolumn{2}{c}{$r_{\text{DC}}=10$} \\
             & aciertos & falsos & aciertos & falsos \\
    \midrule
    0{,}2  & 4/4 & 0   & 4/4 & 0   \\
    0{,}1  & 4/4 & 4   & 4/4 & 0   \\
    0{,}05 & 4/4 & 50  & 4/4 & 0   \\
    0{,}02 & 4/4 & 106 & 4/4 & 140 \\
    0{,}01 & 0/4 & 106 & 4/4 & 634 \\
    \bottomrule
  \end{tabular}
  \caption{Picos detectados según la amplitud de la interferencia y el radio excluido alrededor
    del DC (\texttt{umbral} $=0{,}3$; tabla \texttt{ruido\_sensibilidad.csv}).}
  \label{tab:sensibilidad}
\end{table}

\figura{ruido_deteccion.png}{0.75}{Imagen \texttt{camera} limpia, con dos interferencias y su
  espectro centrado. Los círculos rojos son los 4 picos detectados (dos pares conjugados).}%
  {fig:ruido-deteccion}

En el espectro (Figura~\ref{fig:ruido-deteccion}) la interferencia aparece como puntos aislados,
mientras que el contenido se reparte en la mancha central y en líneas, que vienen de los bordes
rectos de la foto y de la cruz que forman los bordes de la imagen (la fuga de la ventana
rectangular, ahora en 2D). Por eso se puede borrar la interferencia sin tocar el resto.